% ==========================================
% 5. Results
% ==========================================
%
% Status: individual. Each author writes the subsection of their own
% experiment, from experiments/experiment_XXX/ (result.csv, observations.md,
% figures/).
%
% Principle: this section reports what was measured, objectively. The
% interpretation (why, what it means, how it compares with other work) belongs to
% Section~\ref{sec:discussion}. We separate observation from interpretation.
%
% Rules for every experiment:
%   - We report all results, including negative, unexpected or failed ones.
%   - Every number has its unit and, when applicable, its variability
%     (standard deviation or confidence interval, over N runs).
%   - Every figure and table has a caption, a label (fig:exp001:name,
%     tab:exp001:name), and is cited in the text.
%   - Graph axes are named and carry units; vector formats (PDF, SVG) are
%     preferred; files are named exp001_name.pdf.
%   - Figures come from the data in experiments/, generated by a script, never
%     edited by hand.
% Example of a figure:
%   \begin{figure}[htbp]
%       \centering
%       \includegraphics[width=0.8\textwidth]{graphs/exp001_name.pdf}
%       \caption{Short, self-sufficient description (what, conditions, units).}
%       \label{fig:exp001:name}
%   \end{figure}

\section{Results}
\label{sec:results}

% ------------------------------------------
% Short framing paragraph
% ------------------------------------------
% TODO [LAST]: We introduce the organization of the section: one subsection per
%              experiment, followed by a factual summary across experiments.

\subsection{EXP-001: [title]}
\label{sec:results-exp001}
% TODO [DURING]: We recall in one sentence the objective and the research
%                question addressed, with a reference to Section~\ref{sec:method-exp001}.
% TODO [DURING]: We present the main results (table and/or figure) and describe
%                them factually: values, trends, variability.
% TODO [DURING]: We report anomalies, failed runs and deviations from the
%                planned protocol, and state their impact on the results.
% TODO [DURING]: We close with a short factual statement of what the data show
%                (no interpretation). (Owner: ...)

\subsection{EXP-002: [title]}
\label{sec:results-exp002}
% TODO [DURING]: Same structure as EXP-001. (Owner: ...)

\subsection{EXP-003: [title]}
\label{sec:results-exp003}
% TODO [DURING]: Same structure as EXP-001. (Owner: ...)

\subsection{EXP-004: [title]}
\label{sec:results-exp004}
% TODO [DURING]: Same structure as EXP-001. (Owner: ...)

\subsection{EXP-005: [title]}
\label{sec:results-exp005}
% TODO [DURING]: Same structure as EXP-001. (Owner: ...)

\subsection{Summary of Results}
\label{sec:results-summary}
% TODO [LAST]: We give one summary table of the key results of all experiments,
%              organized by research question (RQ, experiment, key metric,
%              value). The table is factual; the reading of it is in the discussion.
% TODO [CHECK]: Every number in this section matches result.csv in the
%               corresponding experiment directory.
% TODO [CHECK]: Every experiment announced in Section~\ref{sec:methodology} has
%               its results here, and nothing is reported without its method.
